% Generated from current MATLAB exports. Include in the correct main/SI sequence.
\begingroup\small\setlength{\tabcolsep}{3pt}
\begin{longtable}{>{\raggedright\arraybackslash}p{0.1700\linewidth}>{\raggedright\arraybackslash}p{0.1320\linewidth}>{\raggedleft\arraybackslash}p{0.1320\linewidth}>{\raggedleft\arraybackslash}p{0.1320\linewidth}>{\raggedleft\arraybackslash}p{0.1320\linewidth}>{\raggedleft\arraybackslash}p{0.1320\linewidth}}
\caption{Held-out peak-origin forecast RMSE and coefficient of determination for current canonical runs.} \label{tab:oos-rmse-r2} \\
\toprule
Country & Outcome & RMSE EnKF & RMSE EnKPF & R2 EnKF & R2 EnKPF \\
\midrule
\endfirsthead
\caption[]{Held-out peak-origin forecast RMSE and coefficient of determination for current canonical runs.} \\
\toprule
Country & Outcome & RMSE EnKF & RMSE EnKPF & R2 EnKF & R2 EnKPF \\
\midrule
\endhead
\midrule
\multicolumn{6}{r}{Continued on next page} \\
\midrule
\endfoot
\bottomrule
\endlastfoot
Korea & Q & 421.7074 & 138.9891 & -9.3168 & -0.1207 \\
Korea & R & 994.1157 & 251.6026 & -0.7813 & 0.8859 \\
Japan & Q & 2.74e+06 & 693.3843 & -989,652.2778 & 0.9365 \\
Japan & R & 4.38e+06 & 1,780.9521 & -1.73e+06 & 0.7139 \\
Pakistan & Q & 2.3e+06 & 11,704.197 & -6,611.4301 & 0.8281 \\
Pakistan & R & 5.89e+06 & 36,367.5041 & -13,680.7961 & 0.4783 \\
\end{longtable}
\par\noindent\footnotesize\textit{Note.} All three completed original peak forecasts are retained. Pakistan has a budget-limited initializer; Korea/Japan forms the separate convergence subset. No forecast row is selected by its held-out error. Intervals are latent ensemble quantiles; the residual score is descriptive.
\endgroup
